% Options for packages loaded elsewhere
\PassOptionsToPackage{unicode}{hyperref}
\PassOptionsToPackage{hyphens}{url}
\documentclass[
  a4paper,
]{ltjsarticle}
\usepackage{xcolor}
\usepackage[margin=25mm]{geometry}
\usepackage{amsmath,amssymb}
\setcounter{secnumdepth}{-\maxdimen} % remove section numbering
\usepackage{iftex}
\ifPDFTeX
  \usepackage[T1]{fontenc}
  \usepackage[utf8]{inputenc}
  \usepackage{textcomp} % provide euro and other symbols
\else % if luatex or xetex
  \usepackage{unicode-math} % this also loads fontspec
  \defaultfontfeatures{Scale=MatchLowercase}
  \defaultfontfeatures[\rmfamily]{Ligatures=TeX,Scale=1}
\fi
\usepackage{lmodern}
\ifPDFTeX\else
  % xetex/luatex font selection
\fi
% Use upquote if available, for straight quotes in verbatim environments
\IfFileExists{upquote.sty}{\usepackage{upquote}}{}
\IfFileExists{microtype.sty}{% use microtype if available
  \usepackage[]{microtype}
  \UseMicrotypeSet[protrusion]{basicmath} % disable protrusion for tt fonts
}{}
\makeatletter
\@ifundefined{KOMAClassName}{% if non-KOMA class
  \IfFileExists{parskip.sty}{%
    \usepackage{parskip}
  }{% else
    \setlength{\parindent}{0pt}
    \setlength{\parskip}{6pt plus 2pt minus 1pt}}
}{% if KOMA class
  \KOMAoptions{parskip=half}}
\makeatother
\usepackage{longtable,booktabs,array}
\newcounter{none} % for unnumbered tables
\usepackage{calc} % for calculating minipage widths
% Correct order of tables after \paragraph or \subparagraph
\usepackage{etoolbox}
\makeatletter
\patchcmd\longtable{\par}{\if@noskipsec\mbox{}\fi\par}{}{}
\makeatother
% Allow footnotes in longtable head/foot
\IfFileExists{footnotehyper.sty}{\usepackage{footnotehyper}}{\usepackage{footnote}}
\makesavenoteenv{longtable}
\usepackage{graphicx}
\makeatletter
\newsavebox\pandoc@box
\newcommand*\pandocbounded[1]{% scales image to fit in text height/width
  \sbox\pandoc@box{#1}%
  \Gscale@div\@tempa{\textheight}{\dimexpr\ht\pandoc@box+\dp\pandoc@box\relax}%
  \Gscale@div\@tempb{\linewidth}{\wd\pandoc@box}%
  \ifdim\@tempb\p@<\@tempa\p@\let\@tempa\@tempb\fi% select the smaller of both
  \ifdim\@tempa\p@<\p@\scalebox{\@tempa}{\usebox\pandoc@box}%
  \else\usebox{\pandoc@box}%
  \fi%
}
% Set default figure placement to htbp
\def\fps@figure{htbp}
\makeatother
\setlength{\emergencystretch}{3em} % prevent overfull lines
\providecommand{\tightlist}{%
  \setlength{\itemsep}{0pt}\setlength{\parskip}{0pt}}
% Glyph map for lualatex (newunicodechar). Maps symbols that may lack font glyphs.
\usepackage{newunicodechar}
\usepackage{amssymb}
\newunicodechar{∎}{\ensuremath{\blacksquare}}
\newunicodechar{✓}{\ensuremath{\checkmark}}
\newunicodechar{✗}{\ensuremath{\times}}
\newunicodechar{⟹}{\ensuremath{\Longrightarrow}}
\newunicodechar{⟺}{\ensuremath{\Longleftrightarrow}}
\newunicodechar{↪}{\ensuremath{\hookrightarrow}}
\newunicodechar{⌊}{\ensuremath{\lfloor}}
\newunicodechar{⌋}{\ensuremath{\rfloor}}
\newunicodechar{≃}{\ensuremath{\simeq}}
\newunicodechar{≈}{\ensuremath{\approx}}
\newunicodechar{≤}{\ensuremath{\leq}}
\newunicodechar{≥}{\ensuremath{\geq}}
\newunicodechar{≠}{\ensuremath{\neq}}
\newunicodechar{→}{\ensuremath{\rightarrow}}
\newunicodechar{←}{\ensuremath{\leftarrow}}
\newunicodechar{↔}{\ensuremath{\leftrightarrow}}
\newunicodechar{⊥}{\ensuremath{\perp}}
\newunicodechar{√}{\ensuremath{\surd}}
\newunicodechar{∑}{\ensuremath{\sum}}
\newunicodechar{∏}{\ensuremath{\prod}}
\newunicodechar{∞}{\ensuremath{\infty}}
\newunicodechar{∈}{\ensuremath{\in}}
\newunicodechar{∀}{\ensuremath{\forall}}
\newunicodechar{∣}{\ensuremath{\mid}}
\newunicodechar{γ}{\ensuremath{\gamma}}
\newunicodechar{λ}{\ensuremath{\lambda}}
\newunicodechar{μ}{\ensuremath{\mu}}
\newunicodechar{ρ}{\ensuremath{\rho}}
\newunicodechar{σ}{\ensuremath{\sigma}}
\newunicodechar{φ}{\ensuremath{\varphi}}
\newunicodechar{ψ}{\ensuremath{\psi}}
\newunicodechar{θ}{\ensuremath{\theta}}
\newunicodechar{Δ}{\ensuremath{\Delta}}
\newunicodechar{Π}{\ensuremath{\Pi}}
\newunicodechar{Λ}{\ensuremath{\Lambda}}
\newunicodechar{τ}{\ensuremath{\tau}}
\newunicodechar{ε}{\ensuremath{\varepsilon}}
\newunicodechar{ω}{\ensuremath{\omega}}
\newunicodechar{π}{\ensuremath{\pi}}
\newunicodechar{ν}{\ensuremath{\nu}}
\newunicodechar{±}{\ensuremath{\pm}}
\newunicodechar{×}{\ensuremath{\times}}
\newunicodechar{·}{\ensuremath{\cdot}}
\usepackage{bookmark}
\IfFileExists{xurl.sty}{\usepackage{xurl}}{} % add URL line breaks if available
\urlstyle{same}
\hypersetup{
  pdftitle={The Equation of Motion of a Single Wave in a Relational-Wave Closed System --- Kuramoto-Type Phase Coupling on L(K\_N) and the Condition for a Self-Consistent Floor on the 90-Degree Phase-Difference Skeleton},
  pdfauthor={Noriaki Kihara (WF System Co., Ltd.), ORCID 0009-0004-6753-4020},
  hidelinks,
  pdfcreator={LaTeX via pandoc}}

\title{The Equation of Motion of a Single Wave in a Relational-Wave
Closed System --- Kuramoto-Type Phase Coupling on \(L(K_N)\) and the
Condition for a Self-Consistent Floor on the 90-Degree Phase-Difference
Skeleton}
\author{Noriaki Kihara (WF System Co., Ltd.), ORCID 0009-0004-6753-4020}
\date{2026-09-12}

\begin{document}
\maketitle

\textbf{Series:} Foundations of Self-Consistent Relational-Wave Closed
Systems (Paper 1 of 10)\\
\textbf{Author:} Noriaki Kihara (WF System Co., Ltd.)　\textbf{Date:}
2026-09-12\\
\textbf{Version DOI:} 10.5281/zenodo.22728761\\
\textbf{Concept DOI:} 10.5281/zenodo.22728760\\
\textbf{ORCID:} 0009-0004-6753-4020\\
\textbf{Zenodo:} https://zenodo.org/records/22728761\\

\begin{center}\rule{0.5\linewidth}{0.5pt}\end{center}

\subsection{Abstract}\label{abstract}

Consider the closed-system map that places a complex relational wave
\(z_e\in\mathbb C\) on every edge of the complete graph \(K_N\) and
evolves it by a phase-only real antisymmetric generator (defined in the
prior work {[}1{]}). The question of this paper arises from an empirical
motivation: the floor produced by make\_parent --- which builds a
self-consistent fixed point \(KZ=i\sigma Z\) from a random seed ---
turned out, for all \(N\), to be \textbf{naturally a 90-degree two-axis
phase lattice} (Experiment 4-1, Study 28). We therefore do not ask ``why
90 degrees''; taking 90 degrees as given, we ask \textbf{``can a highly
symmetric floor that rotates self-consistently be constructed on the
90-degree phase-difference skeleton?''} First, decomposing the one-step
map \(z\to\exp(\Delta\tau K)z\) into amplitude and phase, each edge wave
couples in Kuramoto form \(\sin(\varphi_f-\varphi_e)\) only to edges
sharing a vertex, and we derive

\[\frac{dr_e}{d\tau}=\frac12\sum_{f\sim e}r_f\sin\!\big(2(\varphi_f-\varphi_e)\big)\quad(\text{amplitude}),\qquad r_e\frac{d\varphi_e}{d\tau}=\sum_{f\sim e}r_f\sin^2(\varphi_f-\varphi_e)\quad(\text{phase}).\]

On the 90-degree skeleton (all phase differences are multiples of
\(90^\circ\)) the second-harmonic amplitude drive \(\sin(2\psi)\)
\textbf{vanishes termwise}, giving \(\dot r_e=0\) (amplitude freezing).
The remaining self-consistency condition (uniform rotation
\(\dot\varphi_e=\omega\)) reduces to the eigenvalue problem of the
amplitude vector, \(\omega r_e=\sum_{f\sim e}r_f\sin^2\psi_{ef}\)
(i.e.~\(Wr=\omega r\)). Moreover, with
\(D=\operatorname{diag}(e^{i\varphi_e})\), on the 90-degree skeleton the
identity \(\boxed{D^{-1}KD=iW}\) (\(W_{ef}=A_{ef}\sin^2\psi_{ef}\) real
symmetric, non-negative) holds, so the \textbf{complex self-consistency
problem \(Kz=i\omega z\) (\(z=Dr\)) is exactly equivalent to the
eigenvalue problem \(Wr=\omega r\) of a real symmetric non-negative
integer matrix}. If \(W\) is irreducible, Perron--Frobenius guarantees a
unique positive-amplitude principal eigenvector (\(\omega=\rho(W)\)),
and \(z=Dr\) constructs the self-consistent wave on the 90-degree
skeleton (verified irreducible for all \(N=3\)--\(40\)). Its algebraic
exact solution is given by Paper 3 (the make\_parent floor is an
iterative approximation of the Perron eigenvector of the \(W\) fixed by
its own 90-degree labels, a different member from the integer-\(K\)
floor, with different labels and \(W\)). Placed side by side at the same
\(N\), the make\_parent floor has amplitudes that differ across almost
all edges (\(699\) values at \(N=40\)), whereas the intentionally
constructed high-symmetry floor has \(\le 3\) amplitude values and a
markedly higher symmetry (Fig. 2). Orbit reproduction of the exact map
is \(\le3.0\times10^{-15}\); finite-difference verification of the EOM
is \(\le1.4\times10^{-6}\) (amplitude) and \(\le2.4\times10^{-7}\)
(phase); amplitude freezing on the 90-degree skeleton is
\(\le1.1\times10^{-13}\) (\(N=6,22,40\)). Classification: a derived
consequence of the canonical one\_step, with machine-precision
verification.

\begin{center}\rule{0.5\linewidth}{0.5pt}\end{center}

\subsection{1. Introduction}\label{introduction}

This series studies the dynamics of a closed system in which the full
pairwise relations \(M=N(N-1)/2\) among \(N\) entities are complex
relational waves \(z_e\), with the external seed removed. The
\textbf{make\_parent} procedure used initially builds a self-consistent
fixed point (\(KZ=i\sigma Z\)) from a random seed by iteration, and its
floor turned out, for all \(N=3\)--\(40\), to be \textbf{naturally a
90-degree two-axis phase lattice} (relative phase difference
\(90^\circ\); Fig. 1, Experiment 4-1, Study 28). The phase is naturally
\(90^\circ\), but the amplitudes are edge-by-edge scattered (\(699\)
values at \(N=40\), coefficient of variation of a few percent) and the
centroid \(\Sigma z\neq0\), so the symmetry is not high.

\begin{figure}
\centering
\pandocbounded{\includegraphics[keepaspectratio,alt={Fig. 1: Complex plane of the random make\_parent floor (self-consistent fixed point KZ=iσZ from a random seed) for all N=3-40, step 0. For all N every wave lies on two orthogonal axes (relative phase difference 90 degrees) = naturally a 90-degree two-axis lattice. Amplitudes are scattered edge by edge.}]{figs/p1_single_wave_eom_en_fig1.png}}
\caption{Fig. 1: Complex plane of the random make\_parent floor
(self-consistent fixed point KZ=iσZ from a random seed) for all N=3-40,
step 0. For all N every wave lies on two orthogonal axes (relative phase
difference 90 degrees) = naturally a 90-degree two-axis lattice.
Amplitudes are scattered edge by edge.}
\end{figure}

\textbf{Fig. 1} Complex plane of the random make\_parent floor
(\(N=3\)--\(40\), step 0). The self-consistent fixed point built from a
random seed lies, for all \(N\), naturally on two orthogonal axes
(relative phase difference \(90^\circ\)). This is the starting point of
this paper --- the empirical fact that ``90 degrees is given''; the
reason for it (90 degrees is forced by self-consistency plus termwise
freezing) is given by Paper 0.

Taking this empirical fact as the starting point, the question of this
paper is not ``why 90 degrees'' but --- \textbf{taking 90 degrees as
given (it turned out that way naturally)} --- \textbf{``can a highly
symmetric floor that rotates self-consistently be constructed within the
90-degree phase-difference skeleton?''} Indeed, placed side by side at
the same \(N\) (Fig. 2), against the random make\_parent floor
(amplitudes nearly \(M\) distinct values) the intentionally constructed
high-symmetry floors have few amplitude values (high-symmetry v2:
\(\le 2\) values, \(D_N\)-symmetric / integer-\(K\) analytic floor:
\(\le 3\) values, small centroid, relative equilibrium), with markedly
higher symmetry.

\begin{figure}
\centering
\pandocbounded{\includegraphics[keepaspectratio,alt={Fig. 2: Symmetry comparison of floors (same N, phase all 90 degrees). Left = random make\_parent floor (amplitudes nearly M values, centroid ≠ 0); middle = high-symmetry v2 floor (amplitudes ≤ 2 values but largest-class centroid); right = integer-K analytic floor (≤ 3 values, small centroid, relative equilibrium = Paper 3).}]{figs/p1_single_wave_eom_en_fig2.png}}
\caption{Fig. 2: Symmetry comparison of floors (same N, phase all 90
degrees). Left = random make\_parent floor (amplitudes nearly M values,
centroid ≠ 0); middle = high-symmetry v2 floor (amplitudes ≤ 2 values
but largest-class centroid); right = integer-K analytic floor (≤ 3
values, small centroid, relative equilibrium = Paper 3).}
\end{figure}

\textbf{Fig. 2} Symmetry comparison of floors (\(N=6,12,17,40\)).
Against the 90-degree floor generated naturally by random make\_parent
(amplitudes nearly \(M\) values, \(699\) at \(N=40\)), the intentionally
constructed high-symmetry floors have few amplitude values and higher
symmetry (details in Papers 2, 3).

\subsection{2. Dynamics (the map)}\label{dynamics-the-map}

The one\_step of the canonical program
\texttt{run\_N3\_N40\_stage123\_v1.py} (SHA256 \texttt{1abf2353…}) is,
with \(u=e^{i\arg z}\), \(H_{ef}=A_{ef}\,u_f\bar u_e\),
\(H\leftarrow i\,\Im H\): \(z_{n+1}=\exp(-i\Delta\tau H)z_n\).
Rearranged,

\[K_{ef}=A_{ef}\sin(\varphi_f-\varphi_e),\qquad z_{n+1}=\exp(\Delta\tau K)\,z_n,\qquad \Delta\tau=\frac{2\pi}{N},\]

where \(A\) is the adjacency matrix of the line graph \(L(K_N)\) of
\(K_N\) (edges \(e,f\) share a vertex \(\Rightarrow A_{ef}=1\)) and
\(\varphi_e=\arg z_e\). \textbf{\(K\) is determined by phases only and
contains no amplitude; it is a real antisymmetric matrix}, so
\(\exp(\Delta\tau K)\) is a real orthogonal matrix. Acting identically
on \(\Re z\) and \(\Im z\), it exactly conserves \(\|z\|^2\) and the
square-closure \(z^{\mathsf T}z=\sum_e z_e^2\).

\subsection{3. Equation of motion of a single wave
(derivation)}\label{equation-of-motion-of-a-single-wave-derivation}

Writing the generator flow \(dz/d\tau=Kz\) in components, for
\(f\sim e\) (edges sharing a vertex),

\[\frac{dz_e}{d\tau}=\sum_{f\sim e}\sin(\varphi_f-\varphi_e)\,z_f=\sum_{f\sim e}r_f\sin(\varphi_f-\varphi_e)\,e^{i\varphi_f}.\]

With \(z_e=r_e e^{i\varphi_e}\), writing the LHS as
\(e^{i\varphi_e}(\dot r_e+ir_e\dot\varphi_e)\) and using
\(e^{i\varphi_f}=e^{i\varphi_e}e^{i\psi}\)
(\(\psi=\varphi_f-\varphi_e\)) on the RHS to cancel \(e^{i\varphi_e}\)
gives \(\text{RHS}=\sum r_f\sin\psi(\cos\psi+i\sin\psi)\). From the real
and imaginary parts,

\[\boxed{\ \dot r_e=\frac12\sum_{f\sim e}r_f\sin(2\psi)\ }\qquad\boxed{\ r_e\dot\varphi_e=\sum_{f\sim e}r_f\sin^2\psi\ }\]

(using \(\sin\psi\cos\psi=\tfrac12\sin2\psi\)). Each edge wave is an
oscillator coupled in Kuramoto form only to the adjacent edges sharing a
vertex.

\subsection{4. This equation explains the
orbit}\label{this-equation-explains-the-orbit}

\begin{itemize}
\tightlist
\item
  \textbf{Phase (displacement modulo rotation and rigid rotation):}
  since \(\sin^2\ge0\), every wave \(\varphi_e\) advances monotonically
  = rotation. The common part is a rigid rotation; the
  \textbf{difference} of each wave's \((1/r_e)\sum r_f\sin^2\psi\) is
  the phase displacement with the rigid rotation removed.
\item
  \textbf{Amplitude (growth/decay):} the drive is the second harmonic
  \(\sin(2\psi)\) --- a weighted sum of the sine of twice the phase
  difference with adjacent edges.
\item
  \textbf{Floor (90-degree skeleton):} all phase differences are
  multiples of \(90^\circ\)
  \(\Rightarrow\sin(2\psi)=0\Rightarrow\dot r_e=0\) (termwise amplitude
  freezing). Further, in a configuration where the amplitudes satisfy
  \(Wr=\omega r\) (\(W_{ef}=A_{ef}\sin^2\psi_{ef}\)), all waves rotate
  at the same rate \(\dot\varphi_e=\omega\), forming a self-consistent
  relative equilibrium. The existence and uniqueness of this floor is
  treated in §5.
\end{itemize}

The subject of this paper is the \textbf{construction} of the
self-consistent floor on the 90-degree skeleton; that the floor is
unstable (the qualification for ignition) and the dynamics after leaving
the floor (inflation, re-locking) are not treated here (Papers 2, 5, 6).

\subsection{5. Amplitude freezing and the self-consistency condition on
the 90-degree
skeleton}\label{amplitude-freezing-and-the-self-consistency-condition-on-the-90-degree-skeleton}

90 degrees is a given skeleton (the make\_parent self-consistent floor
turned out that way naturally, §1). On top of that, this section shows
that the 90-degree skeleton freezes the amplitudes termwise, and that
the condition for a self-consistent floor becomes an eigenvalue problem
of the amplitudes.

\textbf{Amplitude freezing (termwise):} the amplitude drive is the
second harmonic \(\sin(2\psi)\), whose zeros on \([0,360^\circ)\) are
only \(\psi=0,90,180,270^\circ\) (by enumeration). Hence if all phase
differences are multiples of \(90^\circ\), then \(\sin(2\psi)=0\) for
every pair, i.e.~\(\dot r_e=0\) holds \textbf{term by term} (termwise).
This is a direct consequence of the generator being \(\sin\psi\) (first
harmonic), so that the amplitude drive is the second harmonic; it is why
the 90-degree skeleton is the skeleton on which amplitudes ``can be
placed without being destroyed.'' (Even at general phases, \(\dot r=0\)
can occur if the per-edge sum cancels = cancellation type; this
corresponds to the lock endpoint.)

\textbf{Theorem (real-symmetrization of the 90-degree skeleton).}
Imposing a self-consistent single rotation (\(\dot\varphi_e=\omega\)
uniform) on the 90-degree skeleton, the phase equation becomes
\(\omega r_e=\sum_{f\sim e}r_f\sin^2\psi_{ef}\), i.e.~\(Wr=\omega r\)
(\(W_{ef}=A_{ef}\sin^2\psi_{ef}\)). This is more than an eigenvalue
problem merely ``appearing'' --- it is an \textbf{exact equivalence}:
for \(\psi_{ef}\in\frac{\pi}{2}\mathbb Z\) and
\(D=\operatorname{diag}(e^{i\varphi_e})\), the identity
\(\sin\psi\,e^{i\psi}=\tfrac12\sin2\psi+i\sin^2\psi=i\sin^2\psi\) (⟺
\(\sin2\psi=0\)) gives
\[\boxed{\ D^{-1}KD=iW\ },\qquad W_{ef}=A_{ef}\sin^2\psi_{ef}\in\{0,1\}\ (\text{real symmetric, non-negative}).\]
Therefore
\(\boxed{\,Kz=i\omega z\ (z=Dr,\ r\ \text{real})\iff Wr=\omega r\,}\).
The complex self-consistent eigen-wave problem on the 90-degree skeleton
is exactly equivalent to the eigenvalue problem of the real symmetric
non-negative integer matrix \(W\).

\textbf{Corollary (existence and uniqueness of the positive-amplitude
self-consistent wave).} If \(W\) is irreducible (its associated graph is
connected), then by Perron--Frobenius there exists, up to scale, a
unique \(r\) with \(Wr=\rho(W)r\), \(r_e>0\). Hence \(\boxed{\,z=Dr\,}\)
is a self-consistent single rotation \(Kz=i\rho(W)z\) on the 90-degree
skeleton. That is, the answer to this paper's question ``can a
self-consistent positive-amplitude wave be built on the 90-degree
skeleton?'' is, under irreducibility of \(W\), \textbf{yes (unique up to
scale)}.

\textbf{Verification (all \(N=3\)--\(40\), follow-up).} For both the
integer-\(K\) floor and the make\_parent floor,
\(\|D^{-1}KD-iW\|\le2\times10^{-12}\); the \(\omega=\rho(W)=\sigma\)
(\(iK\) spectral radius) of \(Wr=\omega r\) agrees; \(r>0\); the overlap
between \(r\) and the \(W\) principal eigenvector is \(1.000000\); and
\(W\) is irreducible (connected) for all \(N\). In summary,
\(\boxed{\text{90-degree skeleton}\overset{D^{-1}KD=iW}{\longrightarrow}\text{real non-negative symmetric eigenproblem}\overset{\rm PF}{\longrightarrow}\text{existence/uniqueness of positive-amplitude self-consistent wave}}\).

The \textbf{closed form} of that high-symmetry solution (the
\(N\)-dependence of \(\rho(W)\), the 2/3-value degeneracy of the
amplitudes, and the agreement with \(\sigma_{\max}\) of the integer
\(K\)) is solved concretely by Paper 3. The make\_parent floor
approximates the Perron principal eigenvector of the \(W\)
(\(W_{\rm mp}\)) fixed by its own 90-degree phase labels obtained
through iteration (the sum rule
\((\sum_f r_f\sin^2\psi_{ef})/r_e=\sigma\) holds for all \(N\), Study
31); it is a different member from the integer-\(K\) floor, with
different labels and different \(W\), yet consistent in that each is the
Perron solution of its own \(W\) (\(W_{\rm mp}\neq W_{\rm integerK}\),
Paper 0).

That is, 90 degrees is not a derived necessity but a given skeleton, and
what this paper establishes is that ``on the 90-degree skeleton the
amplitudes freeze termwise, and the existence of a self-consistent floor
reduces to the eigenvalue problem \(Wr=\omega r\) of the amplitudes.''
The qualification for ignition of the floor (instability) is treated in
Paper 2, the exact solution in Paper 3, and the rapid expansion from it
in Papers 5, 6.

\subsection{\texorpdfstring{6. Numerical verification
(\(N=6,22,40\))}{6. Numerical verification (N=6,22,40)}}\label{numerical-verification-n62240}

{\def\LTcaptype{none} % do not increment counter
\begin{longtable}[]{@{}
  >{\raggedright\arraybackslash}p{(\linewidth - 4\tabcolsep) * \real{0.3333}}
  >{\raggedright\arraybackslash}p{(\linewidth - 4\tabcolsep) * \real{0.3333}}
  >{\raggedright\arraybackslash}p{(\linewidth - 4\tabcolsep) * \real{0.3333}}@{}}
\toprule\noalign{}
\begin{minipage}[b]{\linewidth}\raggedright
Test
\end{minipage} & \begin{minipage}[b]{\linewidth}\raggedright
Content
\end{minipage} & \begin{minipage}[b]{\linewidth}\raggedright
Max error
\end{minipage} \\
\midrule\noalign{}
\endhead
\bottomrule\noalign{}
\endlastfoot
T1 & The exact map \(\exp(\Delta\tau K(\varphi_n))z_n\) reproduces the
canonical orbit \(z_{n+1}\) (\(K\) is the generator) &
\(\le3.0\times10^{-15}\) \\
T2 & The amplitude/phase components of the substep
\((\exp(\varepsilon K)z-z)/\varepsilon\) agree with the EOM of §3 &
amplitude \(\le1.4\times10^{-6}\) / phase \(\le2.4\times10^{-7}\) \\
T3 & On the 90-degree lattice of the floor, \(\dot r_e\approx0\)
(amplitude freezing) & \(|\dot r|\le1.1\times10^{-13}\) \\
\end{longtable}
}

(T2 is due to rounding of the \(\varepsilon\)-scale finite difference;
the formula is exact.) Measured 90-degree skeleton vs.~lock endpoint:
the 90-degree skeleton (step 0) has individual terms
\(|\sin2\psi|\sim10^{-12}\text{–}10^{-14}\) (termwise, lattice deviation
\(0.0^\circ\)); the lock endpoint has individual terms \(\sim1.0\) but
per-edge weighted sums \(\sim10^{-7}\text{–}10^{-13}\) (cancellation,
lattice deviation \(\sim45^\circ\)).

\begin{figure}
\centering
\pandocbounded{\includegraphics[keepaspectratio,alt={Fig. 3: Verification of the single-wave EOM (N=6). Left = amplitude, right = phase; measured (vertical) and derived formula (horizontal) agree on y=x.}]{figs/p1_single_wave_eom_en_fig3.png}}
\caption{Fig. 3: Verification of the single-wave EOM (N=6). Left =
amplitude, right = phase; measured (vertical) and derived formula
(horizontal) agree on y=x.}
\end{figure}

\textbf{Fig. 3} Verification of the single-wave equation of motion
(\(N=6\)).

\begin{figure}
\centering
\pandocbounded{\includegraphics[keepaspectratio,alt={Fig. 4: Amplitude freezing (termwise) on the 90-degree skeleton and locking (cancellation). Left = zeros of sin(2ψ) = k·90 degrees; right = measured distributions of the 90-degree skeleton (termwise, individual terms \textasciitilde0) and locking (cancellation, individual terms O(1)) for N=6.}]{figs/p1_single_wave_eom_en_fig4.png}}
\caption{Fig. 4: Amplitude freezing (termwise) on the 90-degree skeleton
and locking (cancellation). Left = zeros of sin(2ψ) = k·90 degrees;
right = measured distributions of the 90-degree skeleton (termwise,
individual terms \textasciitilde0) and locking (cancellation, individual
terms O(1)) for N=6.}
\end{figure}

\textbf{Fig. 4} The 90-degree skeleton (termwise amplitude freezing) and
the lock endpoint (cancellation), measured (\(N=6\)).

\subsection{7. Claim classification / open
items}\label{claim-classification-open-items}

\begin{itemize}
\tightlist
\item
  Classification: a \textbf{derived consequence} of the canonical
  one\_step (plus machine-precision verification). Verdict: retained.
\item
  Assertion (main result of this paper): on the 90-degree skeleton
  \(D^{-1}KD=iW\) holds, and the complex self-consistency problem
  \(Kz=i\omega z\) is \textbf{exactly equivalent} to the eigenvalue
  problem \(Wr=\omega r\) of the real symmetric non-negative integer
  matrix \(W_{ef}=A_{ef}\sin^2\psi_{ef}\). If \(W\) is irreducible,
  Perron--Frobenius makes the positive-amplitude principal eigenvector
  unique (\(\omega=\rho(W)\)) = \textbf{a self-consistent
  positive-amplitude floor can be constructed on the 90-degree skeleton}
  (verified irreducible at machine precision for all \(N=3\)--\(40\)).
  This is the mathematical answer to ``can a self-consistent
  high-symmetry wave be built on the 90-degree skeleton?''
\item
  Open: the overall scale \(r^2\) of the floor (scale is privately
  held), the moduli of relative equilibria (the 90-degree floor is the
  most symmetric member, not the unique solution), the \(N=3\) exception
  (\(\mu/r^2=-3/2\), rank 1), and the general condition for
  irreducibility of \(W\) at an arbitrary 90-degree label configuration.
  These are treated in Papers 2, 3, 5.
\end{itemize}

\subsection{Related work (structural coincidence, not a derivation
source)}\label{related-work-structural-coincidence-not-a-derivation-source}

\begin{itemize}
\tightlist
\item
  Kuramoto model (Kuramoto 1975/1984):
  \(K_{ef}=A_{ef}\sin(\varphi_f-\varphi_e)\) is a Kuramoto-type phase
  coupling on \(L(K_N)\).
\end{itemize}

\subsection{References}\label{references}

\begin{enumerate}
\def\labelenumi{\arabic{enumi}.}
\tightlist
\item
  Noriaki Kihara, ``The Mechanism of Inflationary Rapid Expansion in a
  Self-Consistent Relational-Wave Closed System,'' Concept DOI
  10.5281/zenodo.22112008. (Source of the definition of the dynamics of
  this paper and of \(\tau\neq\) time.)
\item
  Y. Kuramoto, \emph{Chemical Oscillations, Waves, and Turbulence},
  Springer (1984); ``Self-entrainment of a population of coupled
  non-linear oscillators,'' 1975.
\end{enumerate}

\subsection{Reproduction}\label{reproduction}

The verification scripts and data are packaged in full in
\texttt{位相ロック全N確認\_相対平衡\_20260912/単一波運動方程式\_導出と検証\_20260912/}
(\texttt{verify\_single\_wave\_eom\_20260912.py} = T1/T2/T3,
\texttt{verify\_90deg\_necessity\_20260912.py},
\texttt{make\_figures\_eom\_90deg\_20260912.py}, \texttt{REPORT.md},
\texttt{SHA256SUMS.txt}, \texttt{run\_all.sh}). The \textbf{all-\(N\)
verification of the complex↔real-symmetric equivalence \(D^{-1}KD=iW\)
and Perron--Frobenius (§5)} is in
\texttt{位相ロック全N確認\_相対平衡\_20260912/複素実対称同値\_DKD\_iW\_20260912/}
(\texttt{verify\_DKD\_iW\_equivalence\_20260912.py},
\texttt{results/DKD\_iW\_summary.csv}, README/REPORT/SHA/run\_all;
read-only). The symmetry comparison of Fig. 1 is in
\texttt{同/床の対称性比較\_makeparent\_vs\_高対称\_20260912/}. The input
is the existing den=N npz (\(N=6\): 500 step, \(N=22,40\): 2000 step of
this series). The dynamical map is identical to the canonical
\texttt{run\_N3\_N40\_stage123\_v1.py} (SHA256
\texttt{1abf2353fee2e4f56f05e7a6f149fd086885136beb61ab571b48a56b09691567}),
with no change to the physics equations.

\end{document}
